\section{Datasets and Benchmarks}
\label{sec:datasets}

Evaluating future-trajectory predictions requires both a source of real driving
data with ground-truth ego motion and a principled way to score predicted
trajectories. This section describes the two pillars used here: the
\textbf{nuPlan} dataset, which supplies large-scale real driving logs and
defines the closed-loop planning task, and the \textbf{NAVSIM} benchmark, which
provides a reproducible, simulation-aligned scoring protocol built on the
Predictive Driver Model Score family of metrics. Orbis itself is trained and
evaluated partly on nuPlan, so anchoring our trajectory study to the same data
and metrics keeps the results directly comparable.

\subsection{nuPlan}
\label{sec:nuplan}

nuPlan~\cite{nuplan} is the first large-scale, closed-loop, machine-learning
planning benchmark for autonomous driving. It comprises roughly $1500$ hours of
human driving collected in four cities --- Boston, Pittsburgh, Las Vegas, and
Singapore --- chosen for their contrasting driving cultures (dense casino
pick-up/drop-off zones and eight-lane intersections in Las Vegas, double-parking
in Boston, idiosyncratic left-turn precedence in Pittsburgh, and left-hand
traffic in Singapore). Each city ships with a semantic map and an efficient map
query API, and the logs include lidar point clouds, camera images, localization,
and steering inputs. Agent trajectories are auto-labeled by an offline
perception system (PointPillars with CenterPoint detection and non-causal
tracking) to near-human accuracy, and intervals are tagged with scenario types
--- merges, lane changes, protected and unprotected turns, interactions with
pedestrians and cyclists, double-parked vehicles, and so on --- which enables
scenario-specific evaluation.

The conceptual contribution of nuPlan is the distinction between
\emph{prediction} and \emph{planning}, and the corresponding move from
open-loop to closed-loop evaluation. Prediction benchmarks ask for the most
likely future trajectories of \emph{other} agents and score them with
displacement metrics (minADE, minFDE, miss rate). Such metrics are ill-suited
to ego planning for two reasons. First, planning is multi-modal at a high level
--- at an intersection, turning left and turning right can be equally valid ---
so penalizing any single predicted trajectory by its $L_2$ distance to the one
recorded human path is misleading. Second, open-loop scoring ignores
interaction: whether to continue an overtake or merge back depends on how other
agents react, which only a closed loop can capture. nuPlan therefore defines
three tasks of increasing difficulty: \emph{open-loop} (imitate the human
trajectory, no control), \emph{non-reactive closed-loop} (the planner drives a
simulated vehicle while other agents replay their recorded tracks), and
\emph{reactive closed-loop} (other agents are themselves controlled by a
planning model). A controller tracks the planner's output through a kinematic
motion model, and the resulting driven trajectory is scored on traffic-rule
compliance, human-likeness, vehicle dynamics/comfort, and goal achievement.

For our purposes, nuPlan provides the ground-truth future ego trajectories that
the diffusion probe is trained to predict, the scenario diversity needed to
study multi-modal behaviour (turns and intersections in particular), and the
shared evaluation lineage with Orbis, whose nuPlan-turns split is drawn from the
same source.

\subsection{The NAVSIM benchmark}
\label{sec:navsim}

While nuPlan defines the closed-loop ideal, running full interactive simulation
is computationally expensive and, for vision-based end-to-end models, hard to
make photorealistic. NAVSIM~\cite{navsim} addresses this with
\emph{pseudo-simulation}: a paradigm that keeps the scalability and determinism
of open-loop evaluation on real data while recovering much of what closed-loop
testing measures --- robustness to distribution shift and the mitigation of
causal confusion. The key idea is to pre-generate, before any evaluation, a set
of synthetic observations around each real frame using a driving-specialized
variant of 3D Gaussian Splatting, so that the agent is scored not only on the
original recorded observation but also on plausible nearby states it might
encounter.

Concretely, evaluation proceeds in two stages, illustrated conceptually in
Figure~\ref{fig:navsim}. In \textbf{Stage~1}, the agent predicts a $4$-second
trajectory from the original real observation; the trajectory is executed in a
simplified bird's-eye-view simulation with a kinematic bicycle model and an LQR
controller, and a score plus an endpoint are computed. In \textbf{Stage~2}, the
same scoring is repeated on synthetic observations whose start points are
sampled around the human's true endpoint (laterally up to $\pm2$\,m, and
longitudinally over the physically reachable range, giving more samples for
faster scenarios). Each Stage-2 observation comes with a plausible heading and
motion history obtained by matching to real driving data. The Stage-2 scores are
then combined with a Gaussian weighting based on the proximity of each synthetic
start point to the Stage-1 predicted endpoint, so that futures the agent is
actually likely to reach count more heavily:
\begin{equation}
s_{\text{combined}} = s_1\cdot\sum_i \hat{w}^{\,i}\,s_2^{\,i},
\qquad
\hat{w}^{\,i} = \frac{w^i}{\sum_j w^j},
\quad
w^i = \exp\!\Big(-\frac{\lVert x^i - \hat{x}\rVert^2}{2\sigma^2}\Big),
\end{equation}
where $\hat{x}$ is the Stage-1 endpoint and $x^i$ the start point of the $i$-th
synthetic scene. This two-stage design correlates substantially better with
true closed-loop scores ($R^2=0.8$) than the best open-loop metric
($R^2=0.7$), while requiring roughly $6\times$ fewer environment interactions
than full simulation. The public NAVSIM~v2 leaderboard uses a challenging
\texttt{navhard} split of nuPlan-derived scenes (unprotected turns, dense
traffic) to standardize submissions.

\begin{figure}[t]
\centering
\begin{tikzpicture}[
  font=\small,
  box/.style={draw, rounded corners, align=center, inner sep=3pt, minimum height=8mm},
  stage/.style={box, fill=blue!8},
  obs/.style={box, fill=gray!12},
  synth/.style={box, fill=green!10},
  score/.style={box, fill=orange!14},
  ar/.style={-{Stealth[length=2mm]}, semithick},
  >={Stealth[length=2mm]}
]
% Stage 1
\node[obs] (real) at (0,0) {Real\\observation};
\node[stage] (ag1) at (2.6,0) {Agent\\(planner)};
\node[score] (s1) at (5.2,0) {Stage 1:\\BEV sim\\$\to s_1,\hat{x}$};
\draw[ar] (real) -- (ag1);
\draw[ar] (ag1) -- (s1) node[midway, above, font=\scriptsize] {4\,s traj.};

% Stage 2 synthetic observations
\node[synth] (syn1) at (8.0,1.4) {Synthetic obs.\ $x^1$};
\node[synth] (syn2) at (8.0,0.2) {Synthetic obs.\ $x^2$};
\node[font=\footnotesize] at (8.0,-0.75) {$\vdots$};
\node[synth] (syn3) at (8.0,-1.7) {Synthetic obs.\ $x^i$};
%\node[font=\scriptsize] at (8.0,-2.1) {(3D Gaussian Splatting,\\pre-generated)};
\draw[ar] (s1) -- (syn1);
\draw[ar] (s1) -- (syn2);
\draw[ar] (s1) -- (syn3);

\node[score] (s2) at (11.0,0) {Stage 2:\\score $s_2^{\,i}$\\each scene};
\draw[ar] (syn1) -- (s2);
\draw[ar] (syn2) -- (s2);
\draw[ar] (syn3) -- (s2);

\node[box, fill=red!10] (comb) at (13.6,0) {Combined\\EPDMS\\$s_1\!\cdot\!\sum_i\hat{w}^i s_2^i$};
\draw[ar] (s2) -- (comb);
\draw[ar] (s1.south) .. controls (5.2,-3.1) and (13.6,-3.1) .. (comb.south)
  node[pos=0.5, below, font=\scriptsize] {proximity weights $\hat{w}^i$ to endpoint $\hat{x}$};
\end{tikzpicture}
\caption{NAVSIM pseudo-simulation. Stage~1 scores the agent on the real
observation and records the trajectory endpoint $\hat{x}$. Stage~2 repeats the
scoring on pre-generated synthetic observations rendered with 3D Gaussian
Splatting around plausible nearby start points $x^i$. Stage-2 scores are
combined with a Gaussian proximity weight $\hat{w}^i$ that favours futures close
to the Stage-1 endpoint, yielding the combined EPDMS. This recovers
robustness-to-distribution-shift signal without sequential closed-loop
simulation.}
\label{fig:navsim}
\end{figure}

\subsection{PDM and Extended PDM scores}
\label{sec:pdms}

The metric scored at each NAVSIM stage is the Predictive Driver Model Score
(PDMS) and its extension EPDMS, formalized in the Hydra-MDP++
work~\cite{hydra}. Both follow the same template: a product of binary
\emph{penalty} sub-scores, which can zero out the whole score for a serious
violation, multiplied by a weighted average of graded \emph{quality}
sub-scores. The original PDM score is
\begin{equation}
\text{PDMS} =
\underbrace{\text{NC}\cdot\text{DAC}}_{\text{penalties}}
\cdot
\underbrace{\frac{5\cdot\text{TTC} + 2\cdot\text{C} + 5\cdot\text{EP}}{12}}_{\text{weighted average}},
\label{eq:pdms}
\end{equation}
where NC is no at-fault collision, DAC is drivable-area compliance, TTC is
time-to-collision, C is comfort, and EP is ego progress (route progress relative
to a rule-based reference planner). Because NC and DAC enter multiplicatively, a
collision or driving off the road drives the score to zero regardless of how
smooth or progressive the trajectory was.

The Extended PDM Score adds metrics that capture failures the original score
missed --- traffic-light compliance, lane keeping, driving direction, and a
longer-horizon comfort term:
\begin{equation}
\text{EPDMS} =
\underbrace{\text{NC}\cdot\text{DAC}\cdot\text{DDC}\cdot\text{TL}}_{\text{penalties}}
\cdot
\underbrace{\frac{5\cdot\text{TTC} + 2\cdot\text{C} + 5\cdot\text{EP} + 5\cdot\text{LK} + 5\cdot\text{EC}}{22}}_{\text{weighted average}},
\label{eq:epdms}
\end{equation}
with DDC the driving-direction compliance, TL the traffic-light compliance, LK
lane keeping, and EC extended comfort (consistency of acceleration, jerk, and
yaw across consecutive frames). The NAVSIM~v2 variant additionally introduces a
\emph{filtering} mechanism: if a rule violation was also committed by the human
expert in the same scene, the corresponding penalty is waived, which avoids
penalizing contextually justified maneuvers (for example briefly crossing into
the opposite lane to pass a static obstacle) due to label noise. Table~\ref{tab:metrics}
summarizes the sub-scores, their type, and their weights.

\begin{table}[t]
\centering
\caption{Sub-scores of the PDM and Extended PDM scores. Penalty terms are
multiplicative (a value of $0$ zeroes the whole score); quality terms are
combined in a weighted average. PDMS uses the subset $\{$NC, DAC, TTC, C, EP$\}$;
EPDMS adds DDC, TL, LK, and EC.}
\label{tab:metrics}
\begin{tblr}{
  colspec={lllc},
  row{1}={font=\bfseries},
}
\toprule
Sub-score & Meaning & Type & Weight \\
\midrule
NC  & No at-fault collision        & penalty ($\times$)         & -- \\
DAC & Drivable-area compliance     & penalty ($\times$)         & -- \\
DDC & Driving-direction compliance & penalty ($\times$, EPDMS)  & -- \\
TL  & Traffic-light compliance     & penalty ($\times$, EPDMS)  & -- \\
\midrule
TTC & Time-to-collision            & weighted avg.\             & 5 \\
EP  & Ego progress                 & weighted avg.\             & 5 \\
C   & Comfort                      & weighted avg.\             & 2 \\
LK  & Lane keeping                 & weighted avg.\ (EPDMS)     & 5 \\
EC  & Extended comfort             & weighted avg.\ (EPDMS)     & 5 \\
\bottomrule
\end{tblr}
\end{table}

These metrics give the trajectory study its evaluation language. A trajectory
sampled from the diffusion probe can be scored not only by displacement error
against the human path --- a metric known to correlate poorly with real driving
quality --- but by whether it stays on the drivable area, avoids collisions,
makes progress, and respects driving direction and comfort. Combined with the
distribution-level realism-and-coverage view inherited from Orbis, this lets us
ask the questions that matter for a world-model probe: are the predicted
trajectories not only \emph{close} to what a human did, but \emph{drivable},
\emph{safe}, and appropriately \emph{diverse}.
